\section{Charge Is The Loop Count}

Charge begins when traversal becomes accountable. A local segment may carry
direction, phase, residue, or expectation. It may be essential for the
eventual reading. But a segment that has not returned to its reference has
not yet charged the ledger. It has not come home. The grammar therefore
forbids charge from appearing on an open path merely because something has
moved along it.

The reason is not physical modesty but grammatical necessity. A sign is a
relation to an oriented reference. Until the path closes, the reference has
not been asked to receive what the path carried. A transported distinction
may be locally admissible at each step and still have no final orientation
debt. Only when the carried distinction returns can the grammar ask whether
the return identifies, reverses, or leaves a residue. Charge is the tally of
that closed return.

This is why the chapter calls charge the loop count. The count is not the
number of interior points visited. It is not the elapsed time of the walk. It
is not the magnitude of strain. It is the number of completed traversals that
return to the needle with an orientation the grammar can read. A loop that
returns trivially contributes no signed debt. A loop that returns with a sign
the quotient cannot erase contributes a charged traversal. The ledger counts
closure, not wandering.

The seam gives the count its historical shape. One pass closes. The next
pass opens on the residue left by the completed pass. The record appends; it
does not rewrite itself as though the earlier loop had never occurred. A
completed traversal therefore becomes a seed for the next traversal. Charge
is billed at the seam because the seam is where closure becomes residue and
residue becomes the next admissible opening.

This also protects charge from being confused with mere repetition. Repeating
an open segment can accumulate many marks without closing a single loop. The
marks may become a time count, a length count, or a refinement count, but
they are not yet a charge count. Charge requires return. It asks whether the
whole path, after carrying its distinctions through the tower, comes back to
the reference with a signed residue. If it does, the loop has been counted.

The integer character of the count follows from this closure discipline. A
closed traversal either has entered the ledger as a completed loop or it has
not. A half-remembered loop can carry an unfinished phase, a boundary
obstruction, or an unresolved residue, but it is not half a charge in the
grammar's first reading. Fractional descriptions may later appear as smooth
shadows or averaged presentations. The ledger that bills closed traversal
counts whole returns.

The topological integer-count example illustrates this without becoming the
argument. A thread that winds around a boundary is counted by its winding
number; no fractional thread is admitted in the causal record. The example
shows the same obligation in a familiar shape. What matters for the present
chapter is the closure rule: the loop either closes with a counted winding
or it does not enter as that loop.

The Aharonov-Bohm example gives the phase face of the same rule. Along each
local branch the field reading may be silent, yet the recombined loop can
return with a phase shift. The physical example is richer than the chapter
needs. The grammar takes only the form: local silence does not force global
triviality, and a closed comparison can carry a residue that no open segment
could spend as charge.

This is why the positron in the previous chapter could be read only as a
closed-loop holonomy sign. The positive sign did not live as a free object on
an open path. It appeared when a tilted comparison closed and returned with
the opposite orientation. The flat baseline read the native negative sign;
the tilted loop read the positive sign. In both cases, charge was read only
where closure had made orientation accountable.

The count also keeps the electron disciplined. To name the electron by
counting is not to attach a mythic label to a particle. It is to read the
negative charged traversal of the loop over the native baseline. The sign
belongs to the closed comparison; the count belongs to the completed
traversal; the name belongs to the reading in which those two obligations
meet. If the loop does not close, the name has not been earned.

Charge is therefore the lower face of the backward number. It reads how many
times the tower has come home with a signed orientation. It is lower because
it depends first on closure and count. It is not yet mass, because it has not
asked how the response changes after threshold. It is not yet phase, because
it has not asked which quarter-turn carries the orientation. It is not yet
sameness, because it has not touched the needle as an identification face.

The next obligation is strain. Once a loop can be counted, the grammar can
ask how the loop resists being varied. A single closed traversal reads
charge. A second difference of traversal reads how the response bends past
threshold. That bending is the mass face of the same backward number.
